\begin{center}\textbf{Abstract}\end{center}
SafeInvest must choose suitable investments, avoid unnecessary trading, and recognize when automated advice should stop. We develop three horizon-specific portfolios using a convex expected-log-growth problem constrained by tail log drawdown. A separate execution rule moves partially toward the strategic allocation only when the estimated benefit exceeds transaction costs and a safety margin. Suitability is assessed at each customer's actual horizon using out-of-fold loss outcomes, conservative probability bounds, and a doubled-cost stress test. The selected, frozen Plan A assigns customers B and C to Conservative, E to Balanced, and A and D to human review.

All development uses 120 complete scenarios. A single frozen evaluation uses the remaining 60; its saved outputs provide the evidence reported here. At 30 years and a 0.25\% transaction cost, Balanced A delivers a median real annual growth of \BalancedGrowth\%, compared with \EqualGrowth\% for annual equal-weight rebalancing. Its 90th-percentile maximum drawdown is also higher: \BalancedPninety\% versus \EqualPninety\%. Against annual rebalancing to the same strategic allocation, the execution rule reduces mean real fees by \CostSaving\% and trades at \BalancedTrades\% of subsequent annual reviews. A severe scenario nevertheless loses \SevereLoss\% from peak to trough. The evidence supports selective, explainable automation with explicit escalation; it does not establish a guarantee against loss or universal superiority over the baseline.

\section{Decision Problem and Evidence}
The advisor invests in equities, government bonds, and Treasury bills. Its responsibilities are linked: define Conservative, Balanced, and Growth products; match customers to a product or a human advisor; and choose among keeping holdings, rebalancing, changing the current target, and requesting human review. Optimizing returns alone does not address these service decisions.

The supplied data comprise 180 independent synthetic scenarios, each with 30 annual observations of asset returns and inflation. They are based on selected US series from the JST Macrohistory Database, Release 6, as described in the problem statement \cite{problem,jst}. Scenario numbers identify simulated paths, not historical episodes. We retain each complete path as one statistical unit and use the first $H$ periods when evaluating a customer with horizon $H$.

Scenarios 1--120 form the development partition; scenarios 121--180 are reserved for the final evaluation. The latter was executed once on 17 October 2026 after model settings, customer assignments, and comparators were frozen. The final decision retains calibrated Plan A because subsequent development candidates did not satisfy the complete replacement criteria. Neural-network experiments are therefore outside the submitted strategy. Report preparation uses saved summaries and verification records only.

Our design separates three questions: what long-run allocation is appropriate, whether the customer can be served automatically, and whether an immediate trade is worth its cost. This makes the recommendation auditable and permits a justified refusal to automate.
